\section{Introducción}

En la Parte 1 se especificó un modelo de simulación de eventos discretos de la
entrada, la permanencia, la salida y las incidencias del público en el estadio
de San Mamés, con ocho preguntas (P1--P8) y un diseño experimental de escenarios
X0--X7 más una sensibilidad S1. Esta Parte 2 recoge el modelo implementado en
Python 3 con SimPy, su verificación y validación, y las respuestas a las ocho
preguntas.

La memoria se organiza así. La sección~\ref{sec:modelo} describe el modelo tal
como se ha implementado, y la sección~\ref{sec:desviaciones} reúne, con su
motivo, todo lo que difiere de la Parte 1. La sección~\ref{sec:verificacion}
explica cómo se calculan los indicadores, el método estadístico y la campaña de
réplicas, y las comprobaciones que dan confianza en el modelo. La
sección~\ref{sec:resultados} responde a P1--P8 en orden, la
sección~\ref{sec:sensibilidad} analiza la sensibilidad y la
sección~\ref{sec:conclusiones} cierra con las conclusiones y los límites.

\paragraph{Cómo leer las cifras.} Cada réplica es un partido completo. Salvo que
se indique otra cosa, las cifras son la media entre réplicas $\pm$ la
semianchura del intervalo de confianza (IC) del 95~\% (t de Student, una
observación por réplica), con el número de réplicas $n$. Las diferencias entre
escenarios están \emph{emparejadas}: se calculan réplica a réplica con la misma
semilla y $n$ es el número de parejas. Como no hay datos reales de colas ni de
salidas, los resultados se interpretan como comparaciones entre escenarios y no
como predicciones exactas de San Mamés, tal como advertía la Parte 1.
